\chapter{Introdução}

A pandemia de COVID-19, declarada pela Organização Mundial da Saúde em março de 2020, impôs ao setor de alimentação fora do lar uma reconfiguração acelerada de práticas, rotinas e processos. Medidas de distanciamento social adotadas no Brasil incluíram a suspensão de atividades não essenciais, restrições à ocupação de salões e protocolos sanitários mais rígidos \cite{aquino2020medidas,barcellos2022pandemia}. Bares e restaurantes, fortemente dependentes de presença e interação humana, foram obrigados a adotar rapidamente soluções digitais para manter operação e reduzir contato físico \cite{finkler2020impactos,costa2021implicacoes}. Entre essas soluções, os \emph{cardápios digitais} com \textbf{autoatendimento via \emph{QR Code}} --- nos quais o cliente acessa o menu pelo próprio celular, faz o pedido e, eventualmente, paga sem intermediação de um garçom --- emergiram como uma das mudanças mais visíveis e persistentes do período.

Encerrada a fase aguda da pandemia, esse padrão de digitalização não retrocedeu. Estudos com modelos baseados em agentes sugerem que a difusão de plataformas digitais no \emph{food service} desencadeada pela COVID-19 produz efeitos de longo prazo, com mudanças permanentes no comportamento de consumidores e estabelecimentos \cite{giovanini2024pandemia}. Mesmo após o relaxamento das restrições sanitárias, muitos estabelecimentos mantiveram o cardápio digital como interface preferencial, ao passo que outros voltaram, parcial ou totalmente, ao cardápio físico. A coexistência entre essas configurações --- e a tensão entre ganhos operacionais reivindicados pelos gestores e críticas relatadas por clientes --- motiva a investigação proposta neste trabalho.

\section{Tema}
\label{sec:tema}

O tema deste estudo é a \textbf{adoção de cardápios digitais com autoatendimento em bares e restaurantes brasileiros após 2020}, com foco no caso da cidade de Porto Alegre. Cardápios digitais são interfaces \emph{web} ou aplicativos acessados pelo cliente, em geral por meio de um \emph{QR Code} disponibilizado na mesa, que substituem o cardápio impresso e podem, opcionalmente, integrar fluxos de pedido e pagamento \cite{chen2021consumer,santos2024gerador}. Quando essas interfaces incorporam a capacidade de o próprio cliente registrar e enviar o pedido para a cozinha, sem mediação direta de um garçom, configuram um cenário de \textbf{autoatendimento}.

A transformação digital no setor de alimentação fora do lar antecede a pandemia, mas foi por ela severamente acelerada. \citeonline{harms2021effectuation}, ao estudar empreendedores do ramo na Alemanha durante a primeira onda da COVID-19, identificaram caminhos distintos de inovação de modelo de negócio --- alguns baseados em planejamento, outros em ajustes contínuos --- que conduziram à adoção de soluções digitais como resposta à incerteza. No Brasil, \citeonline{costa2021implicacoes} mostraram, por meio de entrevistas com pequenos negócios alimentícios, que a pandemia funcionou como catalisador externo da inserção de tecnologia, ainda que esta não fosse previamente desejada pelos proprietários. A inserção dessa tecnologia no cotidiano de bares e restaurantes coloca em jogo, simultaneamente, dimensões operacionais (custos, integração com cozinha, manutenção do catálogo), de \textbf{experiência do usuário} (clareza visual, fricções de uso, acessibilidade) e jurídicas (tratamento de dados pessoais coletados via formulários eletrônicos).

Porto Alegre constitui um recorte interessante para investigar esse fenômeno por reunir, num espaço urbano relativamente compacto, uma cena gastronômica diversificada (bairros como Cidade Baixa, Moinhos de Vento e Centro Histórico) e estabelecimentos que atravessaram trajetórias distintas durante e após a pandemia. Apesar disso, no contexto brasileiro ainda há poucos trabalhos que articulam, no mesmo estudo, a perspectiva dos estabelecimentos sobre a adoção e a perspectiva dos clientes sobre a experiência resultante.

\section{Delimitação do tema}
\label{sec:delimitacao}

O escopo desta pesquisa é delimitado em quatro eixos.

\textbf{Recorte geográfico.} O estudo concentra-se em estabelecimentos localizados em Porto Alegre, Rio Grande do Sul, com preferência por dois a três bairros que reúnem alta densidade de bares e restaurantes: \emph{Cidade Baixa}, \emph{Moinhos de Vento} e \emph{Centro Histórico}. Estabelecimentos localizados em outros municípios da região metropolitana são excluídos.

\textbf{Recorte temporal.} O período de interesse é de \emph{março de 2020 a 2026}, com foco analítico na inflexão ``antes'' e ``depois'' da pandemia. Estabelecimentos pré-pandemia que adotaram, mantiveram ou alteraram seus formatos de cardápio entram no recorte; estabelecimentos abertos a partir de 2020 também, na medida em que ilustram o cenário ``nascido digital''.

\textbf{Recorte funcional.} Investigam-se apenas \emph{cardápios digitais com autoatendimento} acessados pelo celular do próprio cliente, em geral via \emph{QR Code} apontando para uma página \emph{web}. Não são objeto deste trabalho: sistemas de gestão de cozinha (\emph{Kitchen Display Systems}), integrações com plataformas de \emph{delivery} (iFood, Rappi, Uber Eats), \emph{Enterprise Resource Planning} (ERP) ou totens físicos de autoatendimento --- ainda que esses sistemas possam ser mencionados como contexto.

\textbf{Recorte de sujeitos.} A pesquisa envolve dois grupos de sujeitos complementares: (i) \emph{gestores ou proprietários} de um a três estabelecimentos selecionados como estudos de caso e (ii) \emph{frequentadores} de bares e restaurantes de Porto Alegre, que responderão a um \emph{survey} sem vínculo direto com os estabelecimentos estudados. Essa combinação permite triangular a percepção declarada pelo lado da oferta com a experiência relatada pelo lado da demanda.

\section{Problema}
\label{sec:problema}

A pergunta de pesquisa que orienta este trabalho é formulada como segue:

\begin{quote}
\emph{Como bares e restaurantes de Porto Alegre que adotaram cardápios digitais com autoatendimento desde 2020 percebem os benefícios, riscos e desafios dessa transição, e como seus clientes vivenciam essa interação?}
\end{quote}

A formulação explicita duas perspectivas que costumam ser tratadas separadamente na literatura: a do \emph{estabelecimento}, em geral abordada sob a ótica da adoção tecnológica e da inovação de modelo de negócio \cite{venkatesh2003useracceptance,harms2021effectuation,costa2021implicacoes}, e a do \emph{cliente}, frequentemente analisada sob a ótica da aceitação e do comportamento do consumidor \cite{venkatesh2012consumer,costa2022comportamento,silva2024design}. O problema deste trabalho é, portanto, situado na \textbf{interseção entre essas duas perspectivas}, com ênfase no contexto pós-2020 e no recorte local de Porto Alegre.

\section{Objetivos}
\label{sec:objetivos}

\subsection{Objetivo geral}
\label{sec:obj-geral}

Analisar, por meio de estudo de caso múltiplo articulado a \emph{survey} com frequentadores, como a adoção de cardápios digitais com autoatendimento a partir de 2020 impacta a operação de bares e restaurantes de Porto Alegre e a experiência de seus clientes.

\subsection{Objetivos específicos}
\label{sec:obj-especificos}

\begin{alineas}
	\item caracterizar o estado da arte sobre cardápios digitais, QR menus e self-ordering em food service;
	\item levantar, com gestores de 1 a 3 estabelecimentos, as motivações, custos e desafios técnicos da adoção;
	\item coletar, via questionário, a percepção de frequentadores quanto a usabilidade, confiança e preferências;
	\item confrontar a percepção dos estabelecimentos com a percepção dos clientes, destacando convergências e tensões;
	\item discutir aspectos técnicos transversais: UX em mobile web, LGPD, acessibilidade e dependência de conectividade.
\end{alineas}

\section{Justificativa}
\label{sec:justificativa}

A relevância deste trabalho pode ser apreciada em dois eixos complementares: prático e acadêmico, com ênfase neste último para a Ciência da Computação.

Do ponto de vista \textbf{prático}, o setor de bares e restaurantes foi um dos mais fortemente impactados pela pandemia de COVID-19 no Brasil. \citeonline{finkler2020impactos}, em revisão narrativa publicada no auge da crise, indicaram que o setor de gastronomia precisaria adequar-se a novos critérios sanitários, comportamentais e tecnológicos para sua recuperação. Anos depois, evidências apontam que muitas das mudanças tecnológicas se consolidaram \cite{giovanini2024pandemia}, exigindo dos gestores --- em sua maioria, pequenas e médias empresas --- decisões de digitalização cujo retorno e cujos riscos ainda não estão totalmente compreendidos. Mapear motivações, ganhos e desafios reais a partir de casos brasileiros tem, portanto, valor direto para a tomada de decisão por parte de proprietários, fornecedores de soluções e associações setoriais.

Do ponto de vista \textbf{acadêmico, especificamente para a Ciência da Computação}, o objeto deste trabalho é tecnológico em sua essência: trata-se de uma \emph{web app} consumida em dispositivos móveis, acessada por \emph{QR Code}, sujeita a regras de \textbf{experiência do usuário} \cite{maia2019usabilidade,koivumaki2008perceptions}, a riscos de adoção descritos por modelos clássicos \cite{davis1989perceived,venkatesh2003useracceptance,venkatesh2012consumer} e a exigências legais derivadas da Lei Geral de Proteção de Dados (Lei~13.709/2018) \cite{martins2021principiologia,almeida2022impactos,boscolo2021lgpd,rocha2025abusividade}. Compreender empiricamente como esses elementos se combinam num contexto real de uso pode \textbf{orientar requisitos de software}, padrões de UX e práticas de privacidade aplicáveis a soluções de cardápio digital, contribuindo para a interface entre engenharia de software, interação humano-computador e direito digital no setor de alimentação. Estudos brasileiros recentes em recortes locais distintos --- como \citeonline{silva2024design} em Quixadá/CE, \citeonline{costa2022comportamento} em Brasília/DF e \citeonline{fioreti2026aplicativos} em estabelecimentos do Maranhão --- já investigaram aspectos isolados desse fenômeno; este trabalho contribui ao agregar Porto Alegre como recorte, ao articular as perspectivas de gestores e clientes em um mesmo desenho e ao adotar a UTAUT2 \cite{venkatesh2012consumer} como referência para o instrumento aplicado aos consumidores finais.
